%% Allgemeines -----------------------------------

% Fortlaufende Nummerierung - Entweder 1, 2, 3, 4, ... oder mit einem semantischen Schema
\def \decisionID {20}

% Datum, an dem der Record anefertigt wurde
\def \decisionDate {26.11.2026}

% Wer hat an dem Record mitgewirkt?
\def \decisionParticipants {Quentin Becker, Niklas Käsehagen, OneTutor (Customer Proxy)}

% Der Status in dem sich die Anforderung in diesem Record aktuell befindet.
% Möglichkeiten: Vorgeschlagen, Akzeptiert, Abgelehnt, Abgelöst
\def \decisionStatus {Akzeptiert}

% ID des Records, der diesen Record ablöst, falls zutreffend
\def \decisionSuperseded {}

% ID des Records, den dieser Record ablöst, falls zutreffend
\def \decisionSuperseding {12}

%% Anforderung -----------------------------------

% Titel der Anforderung
\def \requirementTitle {Es soll auf die XML Struktur der .pptx Dateien zugegriffen werden}

% Kontext der Anforderung
\def \requirementContext {Alle Formate der Microsoft Office Suite basieren auf einem zugrundeliegenden XML-Dokument. In diesem werden die Textabschnitte, Graphiken und Bilder gespeichert. Die Office-Programme interpretieren diesen Inhalt dann zu einem Dokument. }

% Konsequenz der Anforderung
\def \requirementConsequences {Similar Slides soll direkt auf die innere Struktur der Präsentation zugreifen und keine rein visuelle Analyse von außen durchführen. Eine Möglichkeit, um auf diese XML-Struktur zuzugreifen ist die Sprache Visual Basic.}